\section{Negative results register}
Every failed or retracted result in this project is preserved here with its
evidence file. We consider this register a deliverable: the negatives are
what make the surviving claims trustworthy.

\subsection{Benchmark-ladder negatives (Azimuth RES$\to$V1 clean split)}
All numbers are mean per-gene Spearman on V1 mouse (9 genes), train RES
human only unless stated; the bar is the same-split Rule Set 2 replication,
$0.569$ (\texttt{results/beat\_azimuth\_v*.json}).
\begin{itemize}
\item \textbf{v1 ridge 0.268, v2 GNN 0.218}: linear and graph baselines far
below the bar; established that the transfer gap is not a model-capacity
issue at the shallow end.
\item \textbf{v3 GBT 0.480, v4 rank-ensemble 0.474}: published-style
gradient boosting on our feature set closes most of the gap but not all;
ensembling ranks did not help.
\item \textbf{v5 deep-stack, tuned, 0.370}: our strongest-looking
architecture collapsed under a paired bootstrap ($p = 1.0$ for leader
$\geq$ challenger). Two flaws were found: leaky stacking (test-fold deep
features in training) and aggressive hyperparameter tuning on the human
assay. Both were corrected in v6.
\item \textbf{v6 leak-free OOF stack 0.490, $p = 0.9996$}: with
out-of-fold deep meta-features and the leader's exact conservative
hyperparameters, the challenger recovered $+0.12$ over v5 yet still loses.
Conclusion: deep meta-features carry human-assay signal that does not
transfer to mouse.
\item \textbf{v7 multi-assay RES+V2 0.547, $p = 0.146$}: adding the Doench
2016 V2 A375 dataset (4{,}194 guides) to the leader's exact features and
hyperparameters made transfer \emph{worse}, not better. Same-species assay
data does not fix cross-species transfer; assay context dominates data
volume.
\item \textbf{Retracted ``break'' 0.792}: an early run appeared to beat the
bar but mixed FC test guides into RES training data (1{,}836 overlapping
30-mers with V1). Retracted in commit \texttt{a22ba74}; the loader now
enforces \texttt{drug != 'nodrug'} and a hermetic test guards the filter.
\end{itemize}

\subsection{Repair-outcome negatives (Tool 7)}
\begin{itemize}
\item \textbf{Microhomology-enumeration model 0.01--0.11}
(\texttt{results/mhdel\_improve.json}): an explicit MH-deletion enumerator
scored near zero against observed MH-deletion rates; sequence-local MH
features alone do not determine deletion usage. Preserved and replaced by
the GBT marginal model.
\item \textbf{Cross-cell-line transfer failure}
(\texttt{results/outcome\_crosscell.json}): our +1-insertion-rate win over
the official Lindel model on held-out Lib-A (0.614 vs 0.557) does not
transfer to Lib-B in the same cell line (0.576 vs 0.646) or across species
(0.533 vs 0.670). The enumerated per-genotype distribution architecture is
the better inductive bias.
\item \textbf{Frameshift gap} (\texttt{results/lindel\_headtohead.json}):
Lindel leads 0.496 to 0.304 on frameshift frequency; our marginal models
miss the deletion-length structure that drives frameshift.
\end{itemize}

\subsection{Model-class negatives}
The GNN (SequenceGCN) underperforms on every supervised task in this
project: 0.218 on transfer (v2), 0.047--0.053 on repair outcomes
(\texttt{results/outcome\_models.json}), 0.154 on microhomology rate. We
keep it in the toolkit because it wins on the multiplex screening task
(pairwise graph structure is its natural input), and we report its
supervised failures rather than silently dropping it.

\subsection{What the negatives changed}
Each negative altered the project: v5's collapse produced the leak-free
stacking protocol; v7 killed the ``more human data'' hypothesis and ended
benchmark-chasing on this split in favor of the paper; the Tool 7 transfer
failure re-scoped the tool's claims from ``competitive with Lindel'' to
``competitive in-distribution only''; the retraction produced the
contamination guard that every later result inherits.
